% !TeX root=../main.tex
\chapter{مقدمه}
%\thispagestyle{empty}

ناوبری خودران در ربات‌های متحرک تنها به یافتن یک مسیر هندسی میان نقطه شروع و هدف محدود نیست، بلکه مجموعه‌ای پیوسته از ادراک محیط، تشخیص هدف و موانع، تصمیم‌گیری حرکتی و اجرای دقیق فرمان‌ها را دربرمی‌گیرد. در یک محیط واقعی، اطلاعات حسگرها کامل و بدون خطا نیست، موانع ممکن است حرکت کنند و شرایط تماس چرخ‌ها با سطح نیز همواره ثابت باقی نمی‌ماند. ازاین‌رو، سامانه ناوبری باید بتواند تصمیم خود را به‌صورت پیوسته اصلاح کند و هم‌زمان معیارهایی مانند رسیدن به هدف، جلوگیری از برخورد، حفظ فاصله ایمن و نرمی حرکت را در نظر بگیرد \cite{Zhu2021,s25113394}.

این مسئله برای یک اسکوتر چهارچرخ محرک اهمیت بیشتری دارد. جرم و لختی سامانه، محدودیت توقف، لغزش چرخ‌ها و تفاوت پاسخ موتورهای چهارگانه می‌توانند باعث شوند فرمان تولیدشده توسط الگوریتم دقیقاً به حرکت مورد انتظار تبدیل نشود. بنابراین، موفقیت سامانه تنها به کیفیت الگوریتم تصمیم‌گیری وابسته نیست و لازم است لایه‌های ادراک، تصمیم‌گیری، ایمنی و کنترل پایین‌سطح به‌صورت هماهنگ طراحی و پیاده‌سازی شوند.

در پژوهش حاضر، یک سامانه ناوبری خودران مبتنی بر الگوریتم کنشگر--منتقد نرم%
\LTRfootnote{Soft Actor-Critic -- SAC} %
 برای اسکوتر چهارچرخ محرک توسعه داده شده است. در این سامانه، ویژگی‌های هدف، وضعیت حرکتی اسکوتر، داده‌های حسگرهای فراصوت و اطلاعات استخراج‌شده از موانع به عامل داده می‌شوند و شبکه حافظه کوتاه‌مدت بلند%
\LTRfootnote{Long Short-Term Memory -- LSTM} %
 برای پردازش دنباله‌ای ویژگی‌های موانع به کار می‌رود. برای افزایش پایداری آموزش در شبیه‌سازی از آموزش معلم--دانشجو و انتقال مطمئن تر به فضای واقعی نیز از یک سازوکار به نام سپر ایمنی استفاده شده است. سامانه ابتدا در محیط شبیه سازی طراحی و آموزش داده شده و سپس اجزای ادراکی، کنترلی و ایمنی آن روی نمونه واقعی اسکوتر پیاده‌سازی و ارزیابی شده‌اند.

\section{زمینه کلی ناوبری خودران}

یک سامانه ناوبری خودران باید داده‌های حسگرها را دریافت کند، وضعیت نسبی ربات، هدف و موانع را تخمین بزند و بر اساس این اطلاعات، فرمان حرکتی مناسب را تولید نماید. در محیط‌های ایستا و دارای نقشه دقیق، می‌توان بخش مهمی از مسیر را پیش از حرکت محاسبه کرد. بااین‌حال، در محیط‌های پویا، موقعیت و جهت حرکت موانع ممکن است تغییر کند و مسیر از پیش تعیین‌شده دیگر ایمن یا قابل اجرا نباشد. در چنین شرایطی، سامانه باید تصمیم‌های محلی و کوتاه‌مدت را بر اساس آخرین مشاهده از محیط اتخاذ کند.

در سامانه واقعی، مشاهده همواره با عدم قطعیت همراه است. دوربین استریو ممکن است در نواحی کم‌بافت یا تحت شرایط نوری نامناسب عمق نادرست تولید کند، آشکارساز دیداری ممکن است هدف یا مانع را در برخی فریم‌ها از دست بدهد و حسگرهای فراصوت نیز تنها بخش محدودی از پیرامون ربات را پوشش می‌دهند. داده‌های انکودر و حسگر اینرسی نیز تحت تأثیر اغتشاش، لغزش و خطای تجمعی قرار دارند. بنابراین، تصمیم‌گیری تنها بر اساس یک اندازه‌گیری لحظه‌ای قابل اعتماد نیست و سامانه باید اطلاعات چند حسگر و روند زمانی مشاهدات را به‌صورت یکپارچه به کار گیرد.

این ویژگی‌ها باعث می‌شوند مسئله ناوبری واقعی به یک فرایند تصمیم‌گیری مارکوف با مشاهده ناقص%
\LTRfootnote{Partially Observable Markov Decision Process -- POMDP} %
 نزدیک باشد. در چنین مسئله‌ای، عامل به حالت کامل محیط دسترسی ندارد و باید بر اساس یک مشاهده تخمینی، تصمیمی تولید کند که پیامد آن در چند گام بعد نیز مناسب باشد. ازاین‌رو، انتخاب ساختار مشاهده و نحوه استخراج اطلاعات زمانی از داده‌های حسگری، بخش مهمی از طراحی سامانه محسوب می‌شود.

\section{یادگیری تقویتی عمیق در ناوبری}

در یادگیری تقویتی، عامل از طریق تعامل با محیط و دریافت پاداش، سیاستی را می‌آموزد که بازده تجمعی را بیشینه کند. در مسئله ناوبری، رفتار مطلوب را می‌توان با ترکیبی از پاداش نزدیک‌شدن به هدف، جریمه برخورد، جریمه نزدیک‌شدن بیش از حد به موانع، محدودکردن تغییرات ناگهانی فرمان و تشویق رسیدن موفق به هدف تعریف کرد. مزیت اصلی این رویکرد، یادگیری مستقیم یک نگاشت غیرخطی میان مشاهده محیط و فرمان حرکتی است؛ بااین‌حال، کیفیت رفتار نهایی به طراحی فضای مشاهده، فضای عمل، تابع پاداش و تنوع سناریوهای آموزشی وابستگی زیادی دارد \cite{Kober2013RLRobotics,Le2024DRLNavigationSurvey}.

ازآنجاکه فرمان حرکتی اسکوتر شامل سرعت خطی و زاویه‌ای پیوسته است، الگوریتم \lr{SAC} به عنوان هسته تصمیم‌گیری انتخاب شده است. این الگوریتم علاوه بر پشتیبانی از فضای عمل پیوسته، با بیشینه‌سازی آنتروپی%
\LTRfootnote{Enthropy} %
، عامل را به حفظ تنوع رفتاری در طول آموزش تشویق می‌کند و معمولاً از نظر کارایی نمونه و پایداری آموزش عملکرد مناسبی دارد \cite{Haarnoja2018SAC}. بااین‌حال، استفاده از یک الگوریتم توانمند به‌تنهایی برای حل مسئله کافی نیست. عامل ممکن است در شرایطی خارج از توزیع داده‌های آموزشی، فرمانی نامطمئن تولید کند یا در اثر خطای حسگر، وضعیت خطر را به‌درستی تشخیص ندهد. به همین دلیل، سیاست یادگرفته‌شده در این پژوهش در کنار یک لایه ایمنی مستقل استفاده می‌شود.

همچنین، اطلاعات مربوط به چند مانع دارای ماهیت ترتیبی و متغیر هستند. تعداد، فاصله، زاویه و سرعت نسبی موانع می‌توانند در طول زمان تغییر کنند. برای استخراج الگوی زمانی این اطلاعات، از ساختار \lr{LSTM} در معماری عامل استفاده شده است. این ساختار به عامل کمک می‌کند تصمیم خود را تنها بر اساس یک تصویر لحظه‌ای از محیط اتخاذ نکند و تغییرات اخیر مشاهده را نیز در نظر بگیرد.

\section{چالش یکپارچه‌سازی ادراک، تصمیم‌گیری و کنترل}

چالش اصلی این پژوهش، طراحی یک الگوریتم منفرد نیست، بلکه تبدیل مجموعه‌ای از زیرسامانه‌های متفاوت به یک زنجیره عملیاتی منسجم است. لایه ادراک باید از تصاویر استریو، حسگرهای فراصوت، انکودرها و حسگر اینرسی%
\LTRfootnote{Inertial Measurement Unit--IMU} %
، اطلاعاتی با قالب ثابت و قابل استفاده برای عامل تولید کند. لایه تصمیم‌گیری باید این مشاهده را به فرمان پیوسته سرعت خطی و زاویه‌ای تبدیل نماید. سپس لایه کنترل پایین‌سطح باید فرمان تولیدشده را با وجود ناحیه مرده، تفاوت موتورهای چهارگانه، تأخیر و لغزش اجرا کند.

یکی از مسائل مهم در این فرایند، شکاف میان شبیه‌سازی و دنیای واقعی است. در شبیه‌ساز، موقعیت هدف و موانع، سرعت‌ها و فاصله‌ها می‌توانند به‌صورت دقیق در اختیار عامل قرار گیرند، درحالی‌که در سامانه عملی همین اطلاعات باید از داده‌های دیداری و حسگری تخمین زده شوند. اغتشاش عمق، حذف موقت تشخیص، نرخ محدود پردازش و تأخیر ارتباط می‌توانند توزیع مشاهدات واقعی را نسبت به داده‌های آموزشی تغییر دهند. سیاستی که تنها با داده دقیق آموزش دیده باشد، ممکن است در چنین شرایطی افت عملکرد شدیدی نشان دهد.

برای کاهش این مشکل، آموزش در دو مرحله معلم و دانشجو انجام می‌شود. عامل معلم با دسترسی به اطلاعات هندسی دقیق شبیه‌ساز، رفتار پایه را یاد می‌گیرد. سپس عامل دانشجو با مشاهده‌ای نزدیک‌تر به شرایط ادراکی واقعی آموزش داده می‌شود و رفتار معلم به‌عنوان لنگر آموزشی مورد استفاده قرار می‌گیرد. علاوه بر این، سپر ایمنی بر اساس فاصله موانع و سطح خطر، فرمان عامل را بررسی می‌کند و در شرایط ضروری سرعت را محدود یا فرمان را اصلاح می‌نماید. این ساختار باعث می‌شود مسئولیت ایمنی تنها بر عهده سیاست یادگرفته‌شده نباشد.

\section{طرح مسئله}

مسئله اصلی پژوهش، طراحی و پیاده‌سازی سامانه‌ای است که اسکوتر چهارچرخ محرک را در محیط دارای موانع ایستا و متحرک، بدون برخورد و با فرمان‌های قابل اجرا به سمت هدف هدایت کند. عامل باید بر اساس مشاهده‌ای شامل موقعیت نسبی هدف، وضعیت حرکتی اسکوتر، داده‌های نزدیک‌برد فراصوت، فرمان قبلی، سطح خطر و ویژگی‌های استخراج‌شده از موانع، سرعت خطی و زاویه‌ای مناسب را تولید نماید.

راهکار پیشنهادی باید چند شرط را به‌طور هم‌زمان برآورده سازد. نخست، عامل باید در محیط شبیه‌سازی توانایی حرکت هدف‌محور و اجتناب از موانع را یاد بگیرد. دوم، جایگزینی داده‌های دقیق شبیه‌ساز با داده‌های ادراکی نباید باعث افت کنترل‌نشده عملکرد شود. سوم، فرمان سیاست باید از نظر ایمنی بررسی شده و سپس به فرمان‌های قابل اجرای موتورهای چهارگانه تبدیل گردد. چهارم، کل زنجیره از حسگر تا موتور باید روی نمونه واقعی اسکوتر قابل پیاده‌سازی و ارزیابی باشد.

بر این اساس، پژوهش حاضر تنها به مقایسه چند الگوریتم یا گزارش منحنی پاداش محدود نیست. موضوع اصلی، ارزیابی عملکرد یک سامانه کامل است که در آن خطای هر زیرسامانه می‌تواند بر نتیجه نهایی اثر بگذارد. بنابراین، علاوه بر ارزیابی عامل در شبیه‌سازی، دقت فاصله‌سنجی، تشخیص هدف و موانع، خروجی حسگر اینرسی، پاسخ موتورهای محرک و عملکرد یکپارچه ناوبری عملی نیز بررسی می‌شوند.

\section{اهمیت و ضرورت تحقیق}

روش‌های کلاسیک ناوبری در محیط‌های ساختارمند و دارای مدل دقیق، عملکرد قابل‌اعتمادی دارند؛ اما در محیط‌های پویا، استفاده از آن‌ها معمولاً به تنظیم دقیق مدل حرکت، تابع هزینه و قواعد اجتناب از برخورد وابسته است. روش‌های یادگیری تقویتی می‌توانند بخشی از این نگاشت پیچیده را از طریق تعامل با محیط بیاموزند، اما در مقابل با چالش‌هایی مانند ایمنی، تعمیم‌پذیری و انتقال از شبیه‌سازی به عمل مواجه‌اند. بررسی هم‌زمان این ظرفیت‌ها و محدودیت‌ها برای یک سامانه واقعی، از نظر علمی و مهندسی اهمیت دارد.

از دیدگاه عملی، موفقیت یک سیاست در شبیه‌سازی به معنای آماده‌بودن آن برای اجرای واقعی نیست. نرخ پردازش تصویر، کیفیت دوربین، تأخیر شبکه، خطای اودومتری%
\LTRfootnote{Odometry} %
، پاسخ موتور و محدودیت سخت‌افزار همگی بر رفتار نهایی اثر می‌گذارند. پیاده‌سازی سامانه روی اسکوتر واقعی امکان شناسایی این گلوگاه‌ها و ارزیابی واقع‌بینانه عملکرد را فراهم می‌کند. افزون بر این، جرم و توان حرکتی اسکوتر ایجاب می‌کند که ایمنی با یک سازوکار مستقل از سیاست نیز پشتیبانی شود.

\section{نوآوری‌های تحقیق}

نوآوری پژوهش حاضر در معرفی یک الگوریتم پایه جدید نیست، بلکه در توسعه و یکپارچه‌سازی اجزای لازم برای ناوبری یادگیری‌محور یک اسکوتر چهارچرخ محرک قرار دارد. مهم‌ترین دستاوردهای پژوهش عبارت‌اند از:

\begin{enumerate}
    \item طراحی یک معماری یکپارچه ادراک--تصمیم--ایمنی--کنترل که از محیط شبیه‌سازی تا نمونه واقعی امتداد دارد
    \item تعریف مشاهده‌ای ساختارمند شامل اطلاعات هدف، وضعیت حرکتی اسکوتر، حسگرهای فراصوت و ویژگی‌های چند مانع
    \item توسعه سیاست \lr{LSTM-SAC} برای تولید هم‌زمان فرمان پیوسته سرعت خطی و زاویه‌ای
    \item به‌کارگیری آموزش معلم--دانشجو برای کاهش وابستگی سیاست به اطلاعات دقیق شبیه‌ساز و آماده‌سازی آن برای داده‌های ادراکی
    \item استفاده از سپر ایمنی برای اصلاح فرمان در شرایط خطرناک و پیاده‌سازی حلقه‌بسته سامانه روی چهار موتور اسکوتر
    \item پیاده‌سازی و ارزیابی عملی زیرسامانه‌های ادراک، تخمین وضعیت، کنترل موتور و ناوبری یکپارچه روی نمونه واقعی
\end{enumerate}

\section{روش کلی انجام تحقیق}

پژوهش با مرور روش‌های کلاسیک و یادگیری‌محور ناوبری و استخراج الزامات مسئله آغاز شده است. سپس مدل اسکوتر، حسگرها، هدف و موانع در محیط شبیه سازی ایجاد شده‌اند. در این مرحله، فضای مشاهده، فضای عمل، تابع پاداش، شرایط پایان اپیزود و سناریوهای آموزشی تعریف شده و آموزش عامل به‌صورت مرحله‌ای انجام گرفته است.

در ادامه، ابتدا عامل معلم با استفاده از اطلاعات دقیق محیط آموزش دیده و پس از دستیابی به رفتار قابل قبول، عامل دانشجو با ورودی‌های نزدیک‌تر به شرایط ادراک واقعی توسعه یافته است. برای پردازش ویژگی‌های موانع از معماری حافظه‌دار استفاده شده و سپر ایمنی نیز در مسیر اعمال فرمان قرار گرفته است. عملکرد عامل از طریق شاخص‌هایی مانند نرخ موفقیت، برخورد، اتمام زمان، فاصله ایمن، نرمی فرمان و میزان مداخله سپر بررسی می‌شود.

به‌طور هم‌زمان، زیرسامانه‌های عملی شامل دریافت تصویر استریو، تصحیح هندسی تصاویر، تولید نقشه اختلاف مکانی، تشخیص هدف، تشخیص و ردیابی موانع، تخمین فاصله، پردازش داده‌های فراصوت، حسگر اینرسی و انکودرها پیاده‌سازی شده‌اند. داده‌های این بخش پس از پالایش و مدیریت مقادیر نامعتبر، به مشاهده ورودی عامل تبدیل می‌شوند.

در معماری عملی، آردوینو%
\LTRfootnote{Arduino} %
 وظیفه دریافت داده‌های پایین‌سطح و کنترل موتورهای چهارگانه، رزبری پای%
 \LTRfootnote{Raspberry Pi} %
  وظیفه تجمیع و انتقال داده‌ها و رایانه پردازشی وظیفه پردازش تصویر و اجرای مدل یادگیری را بر عهده دارد. خروجی عامل پس از عبور از منطق ایمنی به فرمان حرکتی تبدیل و توسط کنترلر حلقه‌بسته موتورها اجرا می‌شود. در پایان، زیرسامانه‌ها به‌صورت جداگانه و سپس سامانه کامل در سناریوهای عملی مورد آزمون قرار می‌گیرند.

\section{ساختار پایان‌نامه}

این پایان‌نامه در هفت فصل تنظیم شده است. فصل دوم به مرور روش‌های کلاسیک ناوبری، روش‌های یادگیری ماشین، یادگیری تقویتی و پژوهش‌های مرتبط با ناوبری ایمن اختصاص دارد و شکاف پژوهشی را مشخص می‌کند. در فصل سوم، مبانی نظری یادگیری تقویتی، فرایند تصمیم‌گیری مارکوف، مشاهده ناقص و مقایسه الگوریتم‌های مناسب برای فضای عمل پیوسته ارائه می‌شوند و دلیل انتخاب الگوریتم \lr{SAC} بیان می‌گردد.

فصل چهارم معماری کلی سامانه پیشنهادی، مدل‌سازی اسکوتر و محیط شبیه‌سازی، لایه ادراک، فضای مشاهده و عمل، تابع پاداش، سپر ایمنی، معماری \lr{LSTM-SAC} و فرایند آموزش معلم--دانشجو را شرح می‌دهد. فصل پنجم به پیاده‌سازی عملی سامانه، آماده‌سازی سخت‌افزار، استقرار حسگرها، اجرای لایه ادراک، کنترل موتورهای چهارگانه و آماده‌سازی آزمون‌های واقعی اختصاص دارد.

در فصل ششم، نتایج آموزش و ارزیابی سیاست‌ها در شبیه‌سازی، عملکرد سپر ایمنی، دقت اجزای ادراکی و کنترلی و نتایج اعتبارسنجی یکپارچه سامانه روی نمونه واقعی اسکوتر ارائه و تحلیل می‌شوند. در نهایت، فصل هفتم به بحث و تحلیل کلی یافته‌ها، نتیجه‌گیری، بیان محدودیت‌های پژوهش و ارائه پیشنهادهایی برای توسعه‌های آینده اختصاص دارد.
